% ============================================================================
% Section: When Does Adaptive Meta-Learning Pay Off?
% ============================================================================
% Purpose: Characterize the cost-benefit frontier across prior quality
% regimes using the prior degradation sweep (N=100 seeds).
% All Corralling results use w₀=0.7, hybrid LinUCB default.
% Updated: February 21, 2026 — post bug-fix, 100-seed results.

% Crossover value from prior degradation experiment (N=100 seeds, w₀=0.7).
\providecommand{\CrossoverWC}{0.55}

\subsection{When Does Adaptive Meta-Learning Pay Off?}
\label{sec:metalearning_cost}

Having established that semantic priors enable strong routing
performance (\S\ref{sec:results}), we now ask a sharper question:
\textbf{when should a practitioner use Corralling (the library's
default) vs tabula rasa?}

The answer---a prior degradation sweep across the full spectrum from
correct to adversarial priors---reveals a crossover at
$\alpha{\approx}0.55$ that separates two regimes: one where Corralling
dominates both alternatives, and one where its variance makes it
the worst choice.

% ----------------------------------------------------------------------------
% Experimental Setup
% ----------------------------------------------------------------------------
\paragraph{Experimental Setup.}
We evaluate three routing strategies across 11~prior corruption levels
on the 750-prompt LMSYS holdout set (2~models, moderate distribution
shift).  Prior corruption interpolates between correct ($\alpha{=}0$)
and model-swapped ($\alpha{=}1$) priors:
%
\begin{equation}
    \mathbf{b}_{\text{new}}^{(i)} = (1{-}\alpha)\,\mathbf{b}_{\text{orig}}^{(i)}
    + \alpha\,\mathbf{b}_{\text{orig}}^{(j)}, \quad i \neq j
\end{equation}
%
At $\alpha{=}0$, priors are correct; at $\alpha{=}0.5$, priors are
uninformative (identical for both models); at $\alpha{=}1.0$, priors
are adversarial (systematically wrong).

All experiments use $N{=}100$ independent seeds, learning rate
$\eta{=}\sqrt{\ln 2/750}{\approx}0.030$, $\gamma{=}0.05$, initial
warmup weight $w_0{=}0.7$, cost penalty $\lambda{=}0$ (quality-only),
warmup $\alpha{=}2.0$ (constant), and tabula rasa
$\alpha{=}1.0 \to 0.01$ (decaying).  All banditGPT conditions use the
production \texttt{BanditRouter} with hybrid LinUCB policy and
family-based parameter sharing.

% ----------------------------------------------------------------------------
% Main Result: Two Regimes
% ----------------------------------------------------------------------------
\paragraph{Two Regimes, One Crossover.}
Figure~\ref{fig:prior_degradation} presents the core finding:
Corralling achieves the lowest regret for good-to-moderate priors
($\alpha \leq 0.5$), substantially beating both warmup-only and
tabula rasa, but degrades sharply beyond $\alpha{\approx}0.55$.

\emph{Good-to-moderate priors ($\alpha \leq 0.5$):}
\begin{itemize}[nosep]
    \item \textbf{Corralling:} 32--38 cumulative regret ($\sigma{\approx}4$)
    \item \textbf{Warmup-only:} 52--59 regret (constant $\alpha$,
          no meta-learning)
    \item \textbf{Tabula Rasa:} ${\sim}52$ regret (prior-independent
          by construction)
\end{itemize}

Corralling's advantage comes from the complementary alpha schedules:
the warmup expert provides informed early selections (constant
$\alpha{=}2.0$), while the tabula rasa expert converges to the best
policy (decaying $\alpha$).  The meta-learner exploits both---fast
initialization \emph{and} eventual exploitation---achieving 30--40\%
lower regret than either alone.

\emph{Adversarial priors ($\alpha \geq 0.6$):}
\begin{itemize}[nosep]
    \item \textbf{Corralling:} 65--78 regret ($\sigma{\approx}23$)
    \item \textbf{Warmup-only:} 59--67 regret (monotonic degradation)
    \item \textbf{Tabula Rasa:} ${\sim}52$ regret (invariant)
\end{itemize}

The warmup expert's constant exploration amplifies systematically
wrong beliefs, and some seeds fail to shift weight to the tabula rasa
expert fast enough.  The bimodal variance ($\sigma{\approx}23$ vs.\
${\approx}4$ in the good-prior regime) confirms that the meta-learner
sometimes recovers and sometimes does not.

Warmup-only degrades monotonically (52 to 67) but is \emph{never}
the best strategy at any corruption level.

% ----------------------------------------------------------------------------
% Decision-Theoretic Framing
% ----------------------------------------------------------------------------
\paragraph{Expected-Value Decision Framework.}
We frame strategy selection as an expected-value calculation:
\begin{itemize}[nosep]
    \item \textbf{Corralling's upside:} saves 17.7 cumulative regret
          (34\% reduction vs.\ tabula rasa) when $\alpha \leq 0.4$
    \item \textbf{Corralling's downside:} costs 23.2 extra regret
          (45\% increase) when $\alpha \geq 0.6$
\end{itemize}

The breakeven occurs when a practitioner believes there is roughly a
\textbf{57\% chance that priors are useful} ($\alpha \leq 0.4$).  Above
this threshold, Corralling has lower expected regret; below it, tabula
rasa is safer.

In practice, most deployments use priors derived from a related domain
(e.g., training on general chat, deploying on customer support),
placing them well within the $\alpha \leq 0.4$ regime.  Truly
adversarial priors---where model preferences are systematically
reversed---require an active distribution mismatch that is rare and
typically known in advance.

% ----------------------------------------------------------------------------
% Stability Mechanisms
% ----------------------------------------------------------------------------
\paragraph{Detecting and Escaping Instability.}
The adversarial regime is \emph{detectable at runtime}.  Two mechanisms
built into the router provide partial protection:
\begin{itemize}[nosep]
    \item \textbf{$\gamma$-floor mixing:} the meta-learner always
          allocates ${\geq}5\%$ of traffic to each expert, ensuring
          the tabula rasa expert continues learning.
    \item \textbf{Loss decay:} stale expert losses are forgotten,
          preventing early mistakes from permanently anchoring the
          meta-learner.
\end{itemize}
For additional robustness, practitioners can monitor the first 50--100
prompts and switch to tabula rasa if the warmup expert's running
accuracy falls below random chance.

% ----------------------------------------------------------------------------
% Connection to library
% ----------------------------------------------------------------------------
\paragraph{Default Recommendation.}
Use Corralling (the library's default) unless there is strong reason
to believe priors are adversarial.  The 57\% breakeven threshold is
easily cleared for most deployments, and early monitoring provides an
escape hatch for the rare case where it is not.

For deployments where prior quality is genuinely unknown and the cost
of instability is unacceptable (e.g., healthcare, finance), tabula
rasa provides a safe, prior-independent baseline at ${\sim}52$ regret.
